\section{Results}
\label{sec:results}

\subsection{Evaluation protocol and matched budgets}
The evaluation compares correctly conditioned, shared-null and cyclically mislabelled models under
matched architecture, training exposure, checkpoint selection and inference budgets. The evaluation
population contains \PendingValue{held-out requests per family per seed} requests per family across
\PendingValue{independent training seeds} independent training seeds. Component-disjoint held-out
products are excluded from training, decoder development and checkpoint selection. Verified exact-L1
yield counts attempts passing molecular validation, decomposition and exact forward replay, with
\emph{all} attempts in the denominator. Post-hoc rejection cannot increase the proposal's per-attempt
acceptance probability (Proposition~\ref{prop:posthoc-budget}). We report means and sample standard
deviations across training seeds, distinguishing retraining variation from conditional Monte Carlo
variation. Raw and final yield count one output per request before and after the declared inference
procedure; failed searches and abstentions remain in the denominator. Appendix~\ref{app:22-results}
specifies the splits, source domains, proposal counts and computational budgets.

\subsection{Shared generation across 22 assembly families}
We first ask whether one conditional flow can cover 22 assembly families while maintaining
performance on the matched three-family comparison. The equal-family exact-L1 yield is
\PendingValue{raw macro L1} for raw readout and \PendingValue{final macro L1} after checked inference
(Table~\ref{tab:22-l1}). The lowest family yield is \PendingValue{lowest family L1}, with
\PendingValue{between-family pattern}. Decomposition coverage, candidate precision and ambiguity
are \PendingValue{decomposition outcomes}; replay of accepted traces is
\PendingValue{accepted-trace replay}. These quantities distinguish an inability to recover a valid
assembly from an inconsistent returned trace. Fixed-coordinate violations and support overflows
occur in \PendingValue{coordinate and support failures} of \PendingValue{total requests} requests.
On common requests for Ugi 3-CR, aza-Michael addition and reductive amination, the change relative to
the three-family model is \PendingValue{three-family bridge effect}
(Table~\ref{tab:22-controls}).

\begin{table}[H]
\caption{\textbf{Exact assembly across 22 ionizable-lipid reaction families.}
Values are mean $\pm$ sample standard deviation across independent training seeds, with all
requests in the yield denominator. Raw and final outputs share request identities; final includes
checked completion, repair and selection. Null and cyclic arms use matched training and inference
budgets. Decomposition coverage and candidate precision are distinct from replay of accepted traces.
Per-seed counts, replay invariants and ambiguity: Appendix~\ref{app:22-results}.}
\label{tab:22-l1}
\centering\footnotesize\setlength{\tabcolsep}{3pt}\renewcommand{\arraystretch}{1.12}
\begin{tabular}{@{}>{\raggedright\arraybackslash}p{.30\textwidth}ccccc@{}}
\toprule
Family & \shortstack{Raw\\L1 (\%)} & \shortstack{Final\\L1 (\%)} & \shortstack{Shared\\null (\%)} & \shortstack{Cyclic\\label (\%)} & \shortstack{Decomp. cov. /\\candidate prec.} \\
\midrule
A3 amine--aldehyde--alkyne & \PendingCell & \PendingCell & \PendingCell & \PendingCell & \PendingCell \\
Acid--epoxide diester & \PendingCell & \PendingCell & \PendingCell & \PendingCell & \PendingCell \\
AEMA & \PendingCell & \PendingCell & \PendingCell & \PendingCell & \PendingCell \\
Aldehyde Ugi 3-CR & \PendingCell & \PendingCell & \PendingCell & \PendingCell & \PendingCell \\
Aldehyde Ugi 4-CR & \PendingCell & \PendingCell & \PendingCell & \PendingCell & \PendingCell \\
$\alpha$-Isocyanoester dihydroimidazole & \PendingCell & \PendingCell & \PendingCell & \PendingCell & \PendingCell \\
Amine alkylation & \PendingCell & \PendingCell & \PendingCell & \PendingCell & \PendingCell \\
Amine--epoxide & \PendingCell & \PendingCell & \PendingCell & \PendingCell & \PendingCell \\
Aryl reductive amination & \PendingCell & \PendingCell & \PendingCell & \PendingCell & \PendingCell \\
Aza-Michael acrylamide & \PendingCell & \PendingCell & \PendingCell & \PendingCell & \PendingCell \\
Aza-Michael acrylate & \PendingCell & \PendingCell & \PendingCell & \PendingCell & \PendingCell \\
Disulfide Michael & \PendingCell & \PendingCell & \PendingCell & \PendingCell & \PendingCell \\
Epoxide O-acylation & \PendingCell & \PendingCell & \PendingCell & \PendingCell & \PendingCell \\
iPhos & \PendingCell & \PendingCell & \PendingCell & \PendingCell & \PendingCell \\
Ketone--isocyanide amide & \PendingCell & \PendingCell & \PendingCell & \PendingCell & \PendingCell \\
Ketone Ugi 4-CR & \PendingCell & \PendingCell & \PendingCell & \PendingCell & \PendingCell \\
Maleate & \PendingCell & \PendingCell & \PendingCell & \PendingCell & \PendingCell \\
O-esterification & \PendingCell & \PendingCell & \PendingCell & \PendingCell & \PendingCell \\
Passerini & \PendingCell & \PendingCell & \PendingCell & \PendingCell & \PendingCell \\
Preassembled thiol-yne & \PendingCell & \PendingCell & \PendingCell & \PendingCell & \PendingCell \\
Reductive amination & \PendingCell & \PendingCell & \PendingCell & \PendingCell & \PendingCell \\
Thiolactone & \PendingCell & \PendingCell & \PendingCell & \PendingCell & \PendingCell \\
\midrule
Equal-family macro & \PendingCell & \PendingCell & \PendingCell & \PendingCell & \PendingCell \\
Pooled all requests & \PendingCell & \PendingCell & \PendingCell & \PendingCell & \PendingCell \\
Worst-family result & \PendingCell & \PendingCell & \PendingCell & \PendingCell & \PendingCell \\
\bottomrule
\end{tabular}
\end{table}

\subsection{Contribution of reaction-program conditioning}
To isolate the effect of reaction information during generation, we compare the conditioned
model with the shared-null and cyclic arms under the same attempt budget and post-hoc verifier.
The conditioned-minus-null difference is \PendingValue{conditioning difference} percentage points,
with \PendingValue{conditioning uncertainty and family variation}. The cyclic control retains
precursor roles, reaction-core coordinates and program depth, isolating the contribution of the
reaction-identity label within a structured program. Its difference from correct conditioning is
\PendingValue{identity-label difference}. The information supplied through layouts, masks, cores and
decoder priors is accounted for separately (Table~\ref{tab:22-controls}); the resulting attribution is
\PendingValue{program versus identity-label conclusion}.

\subsection{Decoder and architecture effects across programs}
With model weights fixed, the raw-to-final exact-L1 change is
\PendingValue{decoder L1 difference}, and the change in outputs satisfying the qualified structural
checks is \PendingValue{decoder design difference} (Table~\ref{tab:22-decoder}). The same source
verifier assesses core-constrained readout, component completion, scoped head retention, role-local
atom/bond proposals and final ring-system repair. Original candidates remain available to selection.
Reselection within the original pool changes yield by \PendingValue{selection-only difference};
adding repaired candidates contributes \PendingValue{repair-pool difference}, at a cost of
\PendingValue{additional proposal and source-check cost}. Larger-pool gains and equal-cost contrasts
are reported separately, because additional search and a changed selection objective do not isolate
the effect of a repair rule.

The architecture comparison separates joint denoising from parameter-matched independent precursor
roles (FACT), cross-attention and gradient-conflict control. Exact-L1 and structural-quality
contrasts are \PendingValue{architecture contrasts}, with
\PendingValue{architecture family dependence} (Table~\ref{tab:22-architecture}). These comparisons
separate learned-model effects from changes made at inference; motif agreement with the same TRAIN
components that define a decoder prior is not independent validation.

\subsection{Precursor novelty, diversity and held-component generalization}
We next ask whether FORGE can preserve exact assembly while generating precursor constitutions
absent from training. Verified component-novel yield is \PendingValue{component-novel yield} per
1,000 requests, including \PendingValue{distinct component-novel yield} distinct products
(Table~\ref{tab:22-diversity}). Finite-inventory methods have zero component novelty relative to
their own visible inventory by construction. Component novelty is verifier-relative and does not
establish product-level catalogue escape. Exact product reachability gives
\PendingValue{catalogue reachability outcome}. Role-specific unique counts and Shannon and Simpson
effective counts are \PendingValue{component diversity outcome}, distinguishing a broad precursor
population from repeated use of a few components.

Held-component recovery is \PendingValue{held-component yield} per 1,000 requests, covering
\PendingValue{held-identity coverage} of reserved identities. Source-disjoint and component-disjoint
exact yields are \PendingValue{generalization yields}, with \PendingValue{size and repeat effects}
across molecular size, closure count and repeated-arm strata (Table~\ref{tab:22-generalization}).
Comparisons with external generators use their qualified chemical scope and a common verifier;
their outcome is \PendingValue{external-baseline comparison}. Unsupported native families remain
unassessed. Fresh noise applied to TRAIN-derived requests is a sampling diagnostic, not evidence
of held-out generalization.

\subsection{Structural quality and agreement with lipid references}
Exact assembly and lipid structural quality measure different properties. The fraction satisfying
exact L1 and all applicable design checks is \PendingValue{raw joint design yield} for raw outputs
and \PendingValue{final joint design yield} after inference
(Table~\ref{tab:22-quality-metrics}). Chemical alerts, loss of a required headgroup site and
ring-condition violations account for \PendingValue{structural failure distribution}. These checks
apply to the complete final graph; source-positive controls give
\PendingValue{control disagreement and abstention rates}. Chemistry-context flags remain separate
from confirmed defects. A retained-basic-site query does not establish pKa, and absence of a narrow
alert does not establish physical stability or delivery activity.

\begin{table}[H]
\caption{\textbf{Structural quality and agreement with design conditions.}
Raw and final arms share requests; values are mean $\pm$ sample standard deviation across training
seeds. Controls are source-positive structures within each metric's qualified scope. Applicable and
assessed counts, control disagreements and abstentions accompany conditional rates; unsupported
checks are not counted as passes.}
\label{tab:22-quality-metrics}\centering\footnotesize
\setlength{\tabcolsep}{3pt}\renewcommand{\arraystretch}{1.12}
\begin{tabular}{@{}>{\raggedright\arraybackslash}p{.62\textwidth}ccc@{}}\toprule
Metric & Raw & Final & Controls\\\midrule
Valid molecules / all requests (\%) & \PendingCell & \PendingCell & \PendingCell \\
Connected molecules / all requests (\%) & \PendingCell & \PendingCell & \PendingCell \\
Exact L1 + all qualified design checks / requests (\%) & \PendingCell & \PendingCell & \PendingCell \\
Molecules with a justified chemical alert / requests (\%) & \PendingCell & \PendingCell & \PendingCell \\
Molecules requiring chemistry-context review / requests (\%) & \PendingCell & \PendingCell & \PendingCell \\
Required head retained / family-scope requests (\%) & \PendingCell & \PendingCell & \PendingCell \\
Correct role-local cycle allocation / requests (\%) & \PendingCell & \PendingCell & \PendingCell \\
Requested ring-size agreement / comparable requests (\%) & \PendingCell & \PendingCell & \PendingCell \\
TRAIN atom-environment coverage (\%) & \PendingCell & \PendingCell & \PendingCell \\
Ring-system TRAIN support / assessed products (\%) & \PendingCell & \PendingCell & \PendingCell \\
Exact L1 + product/component features observed / requests (\%) & \PendingCell & \PendingCell & \PendingCell \\
\bottomrule\end{tabular}
\par\smallskip\raggedright\scriptsize
TRAIN atom/ring support is a reference diagnostic and may share information with decoder priors.
Definitions: Table~\ref{tab:22-metric-definitions}; per-family outcomes:
Tables~\ref{tab:22-structure} and~\ref{tab:22-conditions}.
\end{table}

Agreement with lipid references is \PendingValue{realism outcome} in fingerprint and descriptor
spaces (Table~\ref{tab:22-realism-metrics}). Precision measures proximity to the reference
population, whereas coverage measures how much of that population is represented. Their joint
pattern is \PendingValue{precision-coverage interpretation}. Nearest-reference similarity and
internal diversity distinguish concentrated familiar outputs from diverse but distant structures.
The source-study-grouped two-sample classifier has AUC \PendingValue{two-sample AUC}, with
\PendingValue{two-sample uncertainty}. These comparisons distinguish TRAIN support from
independent lipid-reference agreement; near-chance discrimination does not prove chemical quality.

\begin{table}[H]
\caption{\textbf{Lipid realism and distributional coverage.}
Raw and final outputs are compared within each family and reference population, with
unique-molecule and all-request denominators distinguished. TRAIN-reference diagnostics are
separate from comparisons with source-disjoint observed lipids. Values are mean $\pm$ sample
standard deviation across training seeds.}
\label{tab:22-realism-metrics}\centering\footnotesize
\setlength{\tabcolsep}{3pt}\renewcommand{\arraystretch}{1.12}
\begin{tabular}{@{}>{\raggedright\arraybackslash}p{.49\textwidth}cc>{\raggedright\arraybackslash}p{.24\textwidth}@{}}\toprule
Metric & Raw & Final & Reference\\\midrule
Fingerprint precision (\%) & \PendingCell & \PendingCell & Declared lipid cohort \\
Fingerprint coverage (\%) & \PendingCell & \PendingCell & Declared lipid cohort \\
Fingerprint recall (\%) & \PendingCell & \PendingCell & Declared lipid cohort \\
Distinct reference-supported products / requests (\%) & \PendingCell & \PendingCell & Declared lipid cohort \\
Nearest-reference Tanimoto similarity & \PendingCell & \PendingCell & Declared lipid cohort \\
Internal ECFP4 diversity & \PendingCell & \PendingCell & Generated population \\
Descriptor distribution distance, per feature & \PendingCell & \PendingCell & Declared lipid cohort \\
Descriptor-space precision and coverage & \PendingCell & \PendingCell & Source-disjoint lipids \\
Source-grouped two-sample classifier AUC & \PendingCell & \PendingCell & Source-disjoint lipids \\
\bottomrule\end{tabular}
\par\smallskip\raggedright\scriptsize
ECFP4 uses Morgan radius 2 and 2,048 bits; neighborhood metrics use five neighbors.
Unsupported reference comparisons remain unassessed. Table~\ref{tab:22-descriptors} gives the
full descriptor panel; Table~\ref{tab:22-realism} reports each family and reference population.
\end{table}

\subsection{Recursive synthesis evidence and complete-route yield}
Exact final assembly closes L1 but does not establish an upstream route. Exact precursor-preparation
evidence (L2) covers \PendingValue{L2 component coverage} of required components, and qualified
terminal evidence (L3) covers \PendingValue{current L3 coverage}
(Table~\ref{tab:22-routes}). Complete-product dossier yield is
\PendingValue{complete dossier yield} per 1,000 requests, with
\PendingValue{family-wise evidence depth}. Each complete dossier requires every precursor branch
to terminate in an admitted exact material, preserving role, reaction-step and repeated-component
context. Directly accepted terminals do not require an L2 synthesis. Missing or scope-mismatched
sources, unresolved identities, expired availability and bounded-search censoring are reported
separately. The dominant limitation is \PendingValue{route-closure limitation}; absence of admitted
evidence is not evidence of unsynthesizability. These results concern auditable computational
synthesis dossiers and do not establish experimental synthesis or biological activity.
